\documentclass[a4paper,12pt]{article}
\usepackage[italian]{babel}
\usepackage[utf8]{inputenc}
\usepackage[T1]{fontenc}
\usepackage{geometry}
\usepackage{setspace}
\geometry{margin=2.5cm}
\onehalfspacing

\title{Esercizi sulle Subnet}
\author{Prof. Fedeli Massimo}
\date{}

\begin{document}

\maketitle

\textbf{Data:} \underline{\hspace{3cm}}

\vspace{0.5cm}

Svolgere i seguenti esercizi mostrando, quando richiesto, i passaggi essenziali del procedimento.

\section*{Esercizio 1}
Dato l'indirizzo IP \textbf{192.168.1.34/24}, determina:
\begin{itemize}
    \item indirizzo di rete;
    \item indirizzo di broadcast;
    \item primo host disponibile;
    \item ultimo host disponibile;
    \item numero totale di host utilizzabili.
\end{itemize}

\section*{Esercizio 2}
Considera la rete \textbf{172.16.8.0/23}. Determina:
\begin{itemize}
    \item subnet mask in formato decimale puntato;
    \item indirizzo di broadcast;
    \item intervallo degli host validi;
    \item numero massimo di host disponibili nella rete.
\end{itemize}

\section*{Esercizio 3}
Dato l'indirizzo IP \textbf{10.0.5.129/26}, calcola:
\begin{itemize}
    \item indirizzo di rete;
    \item indirizzo di broadcast;
    \item primo e ultimo host utilizzabili;
    \item dimensione del blocco di indirizzi.
\end{itemize}

\section*{Esercizio 4}
Una rete \textbf{192.168.10.0/24} deve essere suddivisa in \textbf{4 subnet uguali}. Determina:
\begin{itemize}
    \item nuova subnet mask;
    \item numero di host disponibili per subnet;
    \item indirizzi di rete delle quattro subnet.
\end{itemize}

\section*{Esercizio 5}
La rete \textbf{192.168.50.0/24} deve essere suddivisa in \textbf{8 subnet}. Calcola:
\begin{itemize}
    \item nuova subnet mask;
    \item dimensione di ciascuna subnet;
    \item indirizzi di rete delle prime tre subnet;
    \item indirizzo di broadcast della terza subnet.
\end{itemize}

\section*{Esercizio 6}
Dato l'indirizzo \textbf{172.20.10.75/27}, determina:
\begin{itemize}
    \item subnet mask;
    \item indirizzo di rete;
    \item indirizzo di broadcast;
    \item range degli host validi.
\end{itemize}

\section*{Esercizio 7}
Una rete \textbf{10.1.0.0/16} deve supportare \textbf{almeno 500 host per subnet}. Determina:
\begin{itemize}
    \item quanti bit devono essere riservati agli host;
    \item la nuova subnet mask;
    \item il numero massimo di subnet ottenibili;
    \item il numero di host disponibili per subnet.
\end{itemize}

\section*{Esercizio 8}
Dato l'indirizzo \textbf{192.168.100.200/28}, calcola:
\begin{itemize}
    \item dimensione del blocco;
    \item indirizzo di rete;
    \item indirizzo di broadcast;
    \item primo e ultimo host disponibili.
\end{itemize}

\section*{Esercizio 9}
Una rete \textbf{192.168.1.0/24} deve essere suddivisa per tre reparti con le seguenti necessità:
\begin{itemize}
    \item reparto A: 100 host;
    \item reparto B: 50 host;
    \item reparto C: 20 host.
\end{itemize}
Progetta le subnet utilizzando \textbf{VLSM} e indica:
\begin{itemize}
    \item subnet mask di ciascuna rete;
    \item indirizzo di rete;
    \item intervallo host;
    \item indirizzo di broadcast.
\end{itemize}

\section*{Esercizio 10}
Una rete \textbf{172.16.0.0/22} deve essere suddivisa in subnet da \textbf{almeno 60 host}. Determina:
\begin{itemize}
    \item nuova subnet mask;
    \item numero totale di subnet ottenibili;
    \item indirizzo di rete delle prime quattro subnet;
    \item broadcast della quarta subnet.
\end{itemize}

\end{document}
